\documentclass[a4paper,11pt]{article}
\usepackage[T1]{fontenc}
\usepackage[utf8]{inputenc}
\usepackage[left=2cm,right=2cm,top=2.5cm,bottom=2.5cm]{geometry}
\setlength{\headheight}{14pt}
\usepackage{xcolor}
\usepackage{mdframed}
\usepackage{parskip}
\usepackage{fancyhdr}
\usepackage{hyperref}

\definecolor{usercolor}{RGB}{30, 100, 200}
\definecolor{agentcolor}{RGB}{30, 150, 80}
\definecolor{doccolor}{RGB}{200, 90, 10}
\definecolor{userbg}{RGB}{235, 244, 255}
\definecolor{agentbg}{RGB}{235, 250, 238}
\definecolor{docbg}{RGB}{255, 243, 220}
\definecolor{answerbg}{RGB}{250, 250, 235}

\newmdenv[backgroundcolor=userbg,linecolor=usercolor,linewidth=1.5pt,
  leftmargin=0pt,rightmargin=80pt,innerleftmargin=10pt,innerrightmargin=10pt,
  innertopmargin=8pt,innerbottommargin=8pt,roundcorner=5pt,
  skipabove=6pt,skipbelow=3pt]{userbubble}

\newmdenv[backgroundcolor=agentbg,linecolor=agentcolor,linewidth=1.5pt,
  leftmargin=80pt,rightmargin=0pt,innerleftmargin=10pt,innerrightmargin=10pt,
  innertopmargin=8pt,innerbottommargin=8pt,roundcorner=5pt,
  skipabove=3pt,skipbelow=3pt]{agentbubble}

\newmdenv[backgroundcolor=docbg,linecolor=doccolor,linewidth=1pt,
  leftmargin=80pt,rightmargin=0pt,innerleftmargin=10pt,innerrightmargin=10pt,
  innertopmargin=6pt,innerbottommargin=6pt,roundcorner=4pt,
  skipabove=2pt,skipbelow=3pt]{docenv}

\newmdenv[backgroundcolor=answerbg,linecolor=gray,linewidth=0.5pt,
  leftmargin=80pt,rightmargin=0pt,innerleftmargin=10pt,innerrightmargin=10pt,
  innertopmargin=5pt,innerbottommargin=5pt,roundcorner=4pt,
  skipabove=2pt,skipbelow=6pt]{refanswerenv}

\pagestyle{fancy}
\fancyhf{}
\renewcommand{\headrulewidth}{0.3pt}
\cfoot{\thepage}

\lhead{\textcolor{gray}{\small RAG Conversation Dataset}}
\rhead{\textcolor{gray}{\small Conv.\,\#2\ --\ govt0611}}
\begin{document}

\begin{center}
{\LARGE\bfseries RAG Conversation #7}
\par\medskip
{\small\textcolor{gray}{Domain:\ \texttt{mt-rag-govt-elser-512-100-20240611}\quad|\quad Status:\ \textbf{accepted}}}
\end{center}
\vspace{0.4em}
\noindent\textcolor{gray}{\rule{\linewidth}{0.6pt}}
\vspace{0.4em}
\vspace{0.5em}
\noindent{\small\textcolor{usercolor}{\textbf{$\triangleright$\enspace User --- Turn~1}}}\hfill{\tiny\textcolor{gray}{Summarization $\cdot$ ANSWERABLE}}
\begin{userbubble}
What is the goal of NASA’s Europa Clipper mission?
\end{userbubble}
\noindent\hfill{\small\textcolor{agentcolor}{\textbf{Agent\enspace$\triangleleft$}}}
\begin{agentbubble}
The main science goal of NASA's Europa Clipper mission is to determine whether there are places below the surface of Jupiter's icy moon, Europa, that could support life. The mission's three main science objectives are to understand the nature of the ice shell and the ocean beneath it, to investigate its composition, and to characterize its geology. The mission’s detailed exploration of Europa will help scientists better understand the astrobiological potential for habitable worlds beyond our planet.
\end{agentbubble}
\noindent\hfill{\footnotesize\textcolor{doccolor}{\textbf{$\hookrightarrow$\enspace Retrieved documents~(3)}}}
\begin{docenv}
{\footnotesize\textcolor{doccolor}{\textbf{Doc~1}}}\hfill{\tiny\texttt{score:~18.641}}\ {\tiny\textcolor{gray}{\{'relevant': \{'53573702-4140-4812-8658-ae10a54edfc8': \{'value': 'yes', 'timestamp': 1723033284\}\}\}}}\\[3pt]
{\small\textit{Science | Mission – NASA's Europa Clipper}}\\[2pt]
{\tiny\url{https://europa.nasa.gov/mission/science/}}\\[4pt]
{\footnotesize Science | Mission – NASA's Europa Clipper
Science | Mission – NASA's Europa Clipper

NASA

Europa Clipper

Skip Navigation

menu

close modal

Europa Clipper Mission

Science

Mission

MENU

Overview

Science

Timeline

The Team

Overview

History

FAQ

Key Fact

What is Europa Clipper's main science goal? 

Europa Clipper’s main science goal is to determine whether there are places below the surface of Jupiter’s icy moon, Europa, that could support life.

OverviewOverview

Deep beneath the crust of Jupiter’s frozen moon Europa lies a massive liquid water ocean. Exploring this ocean world with our Europa Clipper spacecraft could provide new clues in our search for life beyond Earth. Scheduled to launch in October 2024, Europa Clipper’s main science goal is to determine whether there are places below the surface of Jupiter’s icy moon, Europa, that could support life. Credit: NASA

Jupiter’s moon Europa shows strong evidence for an ocean of liquid water beneath its icy crust. Beyond Earth, Europa is considered one of the most promising currently habitable environments in our solar system. 
Below Europa’s icy surface, evidence suggests there is a global ocean with more water than all of Earth’s oceans combined. Europa could have all the “ingredients” needed for life as we know it: 

Water: Twice as much as Earth’s oceans 

Organics: Essential chemical building blocks from a variety of sources

Energy: Chemical energy sources from the surface and the sea floor

Stability: Conditions remaining similar for 4 billion years

NASA's Europa Clipper spacecraft will perform approximately 50 close flybys of the moon, gathering detailed measurements to investigate whether the moon could have conditions suitable for life. Europa Clipper is not a life detection mission – its main science goal is to determine whether there are places below Europa’s surface that could support life.

Europa Up Close

Take an in-depth look at Jupiter's icy moon and Europa Clipper's plans for investigation.

Learn more ›

Science ObjectivesScience Objectives
Europa Clipper has three main science objectives: 

Determine the Thickness of Europa’s Icy Shell \& How the Ocean Interacts With the Surface}
\end{docenv}
\begin{docenv}
{\footnotesize\textcolor{doccolor}{\textbf{Doc~2}}}\hfill{\tiny\texttt{score:~18.362}}\ {\tiny\textcolor{gray}{\{'relevant': \{'53573702-4140-4812-8658-ae10a54edfc8': \{'value': 'yes', 'timestamp': 1723033420\}\}\}}}\\[3pt]
{\small\textit{New Video Series Captures Team Working on NASA's Europa Clipper – NASA's Europa Clipper}}\\[2pt]
{\tiny\url{https://europa.nasa.gov/news/90/new-video-series-captures-team-working-on-nasas-europa-clipper/}}\\[4pt]
{\footnotesize New Video Series Captures Team Working on NASA's Europa Clipper – NASA's Europa Clipper
The new video series “Spacecraft Makers: Europa Clipper” offers quick updates on the mission’s progress and lifts the curtain on the exacting work that goes into making sure the spacecraft reaches the Jupiter system in 2030. Europa Clipper aims to help answer questions about the ocean that scientists strongly believe lies below Europa’s icy crust.

Learn More

Assembly

Instruments

Meet Europa Clipper

The spacecraft will fly by the moon about 50 times while orbiting Jupiter. (It can’t orbit Europa because doing so would bring Europa Clipper too close to the gas giant’s brutal radiation belts. Learn more in the video.) On each flyby, a suite of science instruments will gather data on how deep the subsurface ocean is, how thick the ice crust is, and, potentially, the characteristics of any plumes that may be venting subsurface water into space. The goal is to find out whether Europa has the potential to support life.
The series’ premiere episode features Europa Clipper Project Manager Jordan Evans, who also has worked on NASA’s Curiosity Mars rover and the agency’s Hubble Space Telescope. In the video, he joins Deputy Science Manager Trina Ray, who worked on NASA’s Cassini and Galileo missions. They venture into JPL’s storied High Bay 1 clean room, where the Europa Clipper spacecraft is under construction – and where all of NASA’s Mars rovers, the twin Voyager spacecraft, and other historic spacecraft were assembled.
While you’re digging into the nuts and bolts of the spacecraft, check out a 24-hour live feed of assembly in progress in High Bay 1.
Additional episodes of “Spacecraft Makers” will include more activity inside other clean rooms where components of Europa Clipper are coming together. Future seasons of the series will cover other missions under construction at JPL.
More About the Mission
Europa Clipper’s main science goal is to determine whether there are places below the surface of Jupiter’s icy moon, Europa, that could support life. The mission’s three main science objectives are to understand the nature of the ice shell and the ocean beneath it, along with the moon’s composition and geology.}
\end{docenv}
\begin{docenv}
{\footnotesize\textcolor{doccolor}{\textbf{Doc~3}}}\hfill{\tiny\texttt{score:~18.014}}\ {\tiny\textcolor{gray}{\{'relevant': \{'53573702-4140-4812-8658-ae10a54edfc8': \{'value': 'yes', 'timestamp': 1723033449\}\}\}}}\\[3pt]
{\small\textit{NASA’s Europa Clipper Survives and Thrives in ‘Outer Space’ on Earth - NASA}}\\[2pt]
{\tiny\url{https://www.nasa.gov/missions/europa-clipper/nasas-europa-clipper-survives-and-thrives-in-outer-space-on-earth/}}\\[4pt]
{\footnotesize NASA’s Europa Clipper Survives and Thrives in ‘Outer Space’ on Earth - NASA
Europa Clipper’s main science goal is to determine whether there are places below the surface of Jupiter’s icy moon, Europa, that could support life. The mission’s three main science objectives are to determine the thickness of the moon’s icy shell and its surface interactions with the ocean below, to investigate its composition, and to characterize its geology. The mission’s detailed exploration of Europa will help scientists better understand the astrobiological potential for habitable worlds beyond our planet.
Managed by Caltech in Pasadena, California, JPL leads the development of the Europa Clipper mission in partnership with the Johns Hopkins Applied Physics Laboratory (APL) in Laurel, Maryland, for NASA’s Science Mission Directorate in Washington. APL designed the main spacecraft body in collaboration with JPL and NASA’s Goddard Space Flight Center in Greenbelt, Maryland. The Planetary Missions Program Office at NASA’s Marshall Space Flight Center in Huntsville, Alabama, executes program management of the Europa Clipper mission.
Find more information about Europa here:
europa.nasa.gov
News Media Contacts
Gretchen McCartneyJet Propulsion Laboratory, Pasadena, Calif.818-393-6215gretchen.p.mccartney@jpl.nasa.gov
Karen Fox / Charles BlueNASA Headquarters, Washington301-286-6284 / 202-802-5345karen.c.fox@nasa.gov / charles.e.blue@nasa.gov
2024-032

Share

Details

Last Updated

Mar 27, 2024

Related TermsEuropa ClipperEuropaJet Propulsion LaboratoryJupiterThe Solar System

 

Explore More

 

6 min read
New Images From Euclid Mission Reveal Wide View of the Dark Universe

Article

	2 days ago	

 

4 min read
NASA’s Psyche Fires Up Its Sci-Fi-Worthy Thrusters

Article

	3 days ago	

 

7 min read
NASA Tool Gets Ready to Image Faraway Planets

Article

	4 days ago	

 

Keep Exploring
Discover Related Topics

Missions

 

Humans in Space

 

Climate Change

 

Solar System

 

The National Aeronautics and Space Administration}
\end{docenv}
\begin{refanswerenv}
{\footnotesize\textit{\textcolor{gray}{\textbf{Reference answer:}\enspace The main science goal of NASA's Europa Clipper mission is to determine whether there are places below the surface of Jupiter's icy moon, Europa, that could support life. The mission's three main science objectives are to understand the nature of the ice shell and the ocean beneath it, to investigate its composition, and to characterize its geology.}}}
\end{refanswerenv}
\vspace{0.5em}
\noindent{\small\textcolor{usercolor}{\textbf{$\triangleright$\enspace User --- Turn~2}}}\hfill{\tiny\textcolor{gray}{Explanation, Factoid $\cdot$ Follow-up $\cdot$ ANSWERABLE}}
\begin{userbubble}
I Europa livable for humans?
\end{userbubble}
\noindent\hfill{\small\textcolor{agentcolor}{\textbf{Agent\enspace$\triangleleft$}}}
\begin{agentbubble}
Europa is not considered directly habitable for humans, but it is one of the most promising environments in our solar system due to the presence of a vast ocean beneath its icy crust. This ocean could potentially contain conditions suitable for life as we know it. The Europa Clipper mission is designed to investigate whether the moon could have conditions suitable for life, but it is not a life detection mission. Europa could have all the necessary "ingredients" for life as we know it: 
- Water: Twice as much as Earth’s oceans
- Organics: Essential chemical building blocks from a variety of sources
- Energy: Chemical energy sources from the surface and the sea floor
- Stability: Conditions remaining similar for 4 billion years
\end{agentbubble}
\noindent\hfill{\footnotesize\textcolor{doccolor}{\textbf{$\hookrightarrow$\enspace Retrieved documents~(3)}}}
\begin{docenv}
{\footnotesize\textcolor{doccolor}{\textbf{Doc~1}}}\hfill{\tiny\texttt{score:~43.121}}\ {\tiny\textcolor{gray}{\{'relevant': \{'53573702-4140-4812-8658-ae10a54edfc8': \{'value': 'no', 'timestamp': 1723033818\}\}\}}}\\[3pt]
{\small\textit{NASA’s Europa Clipper Survives and Thrives in ‘Outer Space’ on Earth - NASA}}\\[2pt]
{\tiny\url{https://www.nasa.gov/missions/europa-clipper/nasas-europa-clipper-survives-and-thrives-in-outer-space-on-earth/}}\\[4pt]
{\footnotesize NASA’s Europa Clipper Survives and Thrives in ‘Outer Space’ on Earth - NASA
Europa Clipper’s main science goal is to determine whether there are places below the surface of Jupiter’s icy moon, Europa, that could support life. The mission’s three main science objectives are to determine the thickness of the moon’s icy shell and its surface interactions with the ocean below, to investigate its composition, and to characterize its geology. The mission’s detailed exploration of Europa will help scientists better understand the astrobiological potential for habitable worlds beyond our planet.
Managed by Caltech in Pasadena, California, JPL leads the development of the Europa Clipper mission in partnership with the Johns Hopkins Applied Physics Laboratory (APL) in Laurel, Maryland, for NASA’s Science Mission Directorate in Washington. APL designed the main spacecraft body in collaboration with JPL and NASA’s Goddard Space Flight Center in Greenbelt, Maryland. The Planetary Missions Program Office at NASA’s Marshall Space Flight Center in Huntsville, Alabama, executes program management of the Europa Clipper mission.
Find more information about Europa here:
europa.nasa.gov
News Media Contacts
Gretchen McCartneyJet Propulsion Laboratory, Pasadena, Calif.818-393-6215gretchen.p.mccartney@jpl.nasa.gov
Karen Fox / Charles BlueNASA Headquarters, Washington301-286-6284 / 202-802-5345karen.c.fox@nasa.gov / charles.e.blue@nasa.gov
2024-032

Share

Details

Last Updated

Mar 27, 2024

Related TermsEuropa ClipperEuropaJet Propulsion LaboratoryJupiterThe Solar System

 

Explore More

 

6 min read
New Images From Euclid Mission Reveal Wide View of the Dark Universe

Article

	2 days ago	

 

4 min read
NASA’s Psyche Fires Up Its Sci-Fi-Worthy Thrusters

Article

	3 days ago	

 

7 min read
NASA Tool Gets Ready to Image Faraway Planets

Article

	4 days ago	

 

Keep Exploring
Discover Related Topics

Missions

 

Humans in Space

 

Climate Change

 

Solar System

 

The National Aeronautics and Space Administration}
\end{docenv}
\begin{docenv}
{\footnotesize\textcolor{doccolor}{\textbf{Doc~2}}}\hfill{\tiny\texttt{score:~42.900}}\ {\tiny\textcolor{gray}{\{'relevant': \{'9dcd142c-415e-4161-a82f-c8b14bfb36ac': \{'value': 'yes', 'timestamp': 1726568616\}\}\}}}\\[3pt]
{\small\textit{Science | Mission – NASA's Europa Clipper}}\\[2pt]
{\tiny\url{https://europa.nasa.gov/mission/science/}}\\[4pt]
{\footnotesize Science | Mission – NASA's Europa Clipper
Science | Mission – NASA's Europa Clipper

NASA

Europa Clipper

Skip Navigation

menu

close modal

Europa Clipper Mission

Science

Mission

MENU

Overview

Science

Timeline

The Team

Overview

History

FAQ

Key Fact

What is Europa Clipper's main science goal? 

Europa Clipper’s main science goal is to determine whether there are places below the surface of Jupiter’s icy moon, Europa, that could support life.

OverviewOverview

Deep beneath the crust of Jupiter’s frozen moon Europa lies a massive liquid water ocean. Exploring this ocean world with our Europa Clipper spacecraft could provide new clues in our search for life beyond Earth. Scheduled to launch in October 2024, Europa Clipper’s main science goal is to determine whether there are places below the surface of Jupiter’s icy moon, Europa, that could support life. Credit: NASA

Jupiter’s moon Europa shows strong evidence for an ocean of liquid water beneath its icy crust. Beyond Earth, Europa is considered one of the most promising currently habitable environments in our solar system. 
Below Europa’s icy surface, evidence suggests there is a global ocean with more water than all of Earth’s oceans combined. Europa could have all the “ingredients” needed for life as we know it: 

Water: Twice as much as Earth’s oceans 

Organics: Essential chemical building blocks from a variety of sources

Energy: Chemical energy sources from the surface and the sea floor

Stability: Conditions remaining similar for 4 billion years

NASA's Europa Clipper spacecraft will perform approximately 50 close flybys of the moon, gathering detailed measurements to investigate whether the moon could have conditions suitable for life. Europa Clipper is not a life detection mission – its main science goal is to determine whether there are places below Europa’s surface that could support life.

Europa Up Close

Take an in-depth look at Jupiter's icy moon and Europa Clipper's plans for investigation.

Learn more ›

Science ObjectivesScience Objectives
Europa Clipper has three main science objectives: 

Determine the Thickness of Europa’s Icy Shell \& How the Ocean Interacts With the Surface}
\end{docenv}
\begin{docenv}
{\footnotesize\textcolor{doccolor}{\textbf{Doc~3}}}\hfill{\tiny\texttt{score:~41.591}}\ {\tiny\textcolor{gray}{\{'relevant': \{'9dcd142c-415e-4161-a82f-c8b14bfb36ac': \{'value': 'yes', 'timestamp': 1726568640\}\}\}}}\\[3pt]
{\small\textit{Europa Clipper - NASA Science}}\\[2pt]
{\tiny\url{https://science.nasa.gov/mission/europa-clipper/}}\\[4pt]
{\footnotesize Europa Clipper - NASA Science
Scientists predict Europa has a salty ocean beneath its icy crust that could hold the building blocks necessary to sustain life.TypeOrbiterLaunchOct. 10, 2024TargetEuropaObjectiveDetermine if Europa has conditions suitable to support lifeMission OverviewSpacecraftWhy Europa?What will Europa Clipper do?Europa Clipper’s main science goal is to determine whether there are places below the surface of Jupiter’s icy moon, Europa, that could support life.The mission’s three main science objectives are to understand the nature of the ice shell and the ocean beneath it, along with the moon’s composition and geology. The mission’s detailed exploration of Europa will help scientists better understand the astrobiological potential for habitable worlds beyond our planet.Read MoreAn artist's concept of NASA’s Europa Clipper spacecraft.NASA/JPL-CaltechWhat is NASA’s Europa Clipper spacecraft?Europa Clipper is a robotic solar-powered spacecraft built to conduct the first detailed investigations of Jupiter's icy moon Europa.With its solar arrays deployed, Europa Clipper spans more than 100 feet (about 30 meters) – about the length of a basketball court. The main body of the spacecraft consists of its avionics vault, radiofrequency module, and propulsion module.Read MoreWorkers use a crane to lift Europa Clipper's high-gain antenna as they prepare to install it on the spacecraft on Aug. 14, 2023. NASA/JPL-CaltechWhy Europa?The search for life beyond Earth is one of NASA’s primary objectives. If humans are to truly understand our place in the universe, we must learn whether our planet is the only place where life exists. So the search is on! There is strong evidence Jupiter's moon Europa has a saltwater ocean that may be one of the best places to look for environments where life could exist beyond Earth.Read MoreThis view of Jupiter’s icy moon Europa was captured by the JunoCam imager aboard NASA’s Juno spacecraft during the mission’s close flyby on Sept. 29, 2022.}
\end{docenv}
\begin{refanswerenv}
{\footnotesize\textit{\textcolor{gray}{\textbf{Reference answer:}\enspace Based on the information provided, Europa is not considered directly livable for humans. However, it is considered one of the most promising currently habitable environments in our solar system due to the presence of a vast ocean beneath its icy crust. This ocean could potentially contain conditions suitable for life as we know it. The Europa Clipper mission is designed to investigate whether the moon could have conditions suitable for life, but it is not a life detection mission.}}}
\end{refanswerenv}
\vspace{0.5em}
\noindent{\small\textcolor{usercolor}{\textbf{$\triangleright$\enspace User --- Turn~3}}}\hfill{\tiny\textcolor{gray}{Keyword $\cdot$ Follow-up $\cdot$ ANSWERABLE}}
\begin{userbubble}
The surface of Jupiter Europa
\end{userbubble}
\noindent\hfill{\small\textcolor{agentcolor}{\textbf{Agent\enspace$\triangleleft$}}}
\begin{agentbubble}
The surface of Jupiter's moon, Europa, is made up of ice. It is covered in a thick crust of frozen water that is estimated to be between 10 to 100 kilometers thick. The surface also shows signs of past and present geological activity, such as ridges, chaos terrain, and lenticulae. These features suggest that there is heat being generated from within Europa, which is thought to be due to tidal heating caused by Jupiter's gravitational pull. Additionally, the surface of Europa is subject to bombardment by radiation from Jupiter, which can cause the ice to darken and form a reddish-brown color in some areas. The surface of Jupiter's moon Europa features a widely varied landscape, including ridges, bands, small rounded domes and disrupted spaces that geologists call "chaos terrain."
\end{agentbubble}
\noindent\hfill{\footnotesize\textcolor{doccolor}{\textbf{$\hookrightarrow$\enspace Retrieved documents~(6)}}}
\begin{docenv}
{\footnotesize\textcolor{doccolor}{\textbf{Doc~1}}}\hfill{\tiny\texttt{score:~48.754}}\ {\tiny\textcolor{gray}{\{'relevant': \{'53573702-4140-4812-8658-ae10a54edfc8': \{'value': 'no', 'timestamp': 1723034206\}\}\}}}\\[3pt]
{\small\textit{Science | Mission – NASA's Europa Clipper}}\\[2pt]
{\tiny\url{https://europa.nasa.gov/mission/science/}}\\[4pt]
{\footnotesize Science | Mission – NASA's Europa Clipper
Science | Mission – NASA's Europa Clipper

NASA

Europa Clipper

Skip Navigation

menu

close modal

Europa Clipper Mission

Science

Mission

MENU

Overview

Science

Timeline

The Team

Overview

History

FAQ

Key Fact

What is Europa Clipper's main science goal? 

Europa Clipper’s main science goal is to determine whether there are places below the surface of Jupiter’s icy moon, Europa, that could support life.

OverviewOverview

Deep beneath the crust of Jupiter’s frozen moon Europa lies a massive liquid water ocean. Exploring this ocean world with our Europa Clipper spacecraft could provide new clues in our search for life beyond Earth. Scheduled to launch in October 2024, Europa Clipper’s main science goal is to determine whether there are places below the surface of Jupiter’s icy moon, Europa, that could support life. Credit: NASA

Jupiter’s moon Europa shows strong evidence for an ocean of liquid water beneath its icy crust. Beyond Earth, Europa is considered one of the most promising currently habitable environments in our solar system. 
Below Europa’s icy surface, evidence suggests there is a global ocean with more water than all of Earth’s oceans combined. Europa could have all the “ingredients” needed for life as we know it: 

Water: Twice as much as Earth’s oceans 

Organics: Essential chemical building blocks from a variety of sources

Energy: Chemical energy sources from the surface and the sea floor

Stability: Conditions remaining similar for 4 billion years

NASA's Europa Clipper spacecraft will perform approximately 50 close flybys of the moon, gathering detailed measurements to investigate whether the moon could have conditions suitable for life. Europa Clipper is not a life detection mission – its main science goal is to determine whether there are places below Europa’s surface that could support life.

Europa Up Close

Take an in-depth look at Jupiter's icy moon and Europa Clipper's plans for investigation.

Learn more ›

Science ObjectivesScience Objectives
Europa Clipper has three main science objectives: 

Determine the Thickness of Europa’s Icy Shell \& How the Ocean Interacts With the Surface}
\end{docenv}
\begin{docenv}
{\footnotesize\textcolor{doccolor}{\textbf{Doc~2}}}\hfill{\tiny\texttt{score:~47.266}}\ {\tiny\textcolor{gray}{\{'relevant': \{'53573702-4140-4812-8658-ae10a54edfc8': \{'value': 'no', 'timestamp': 1723034247\}\}\}}}\\[3pt]
{\small\textit{Europa Clipper - NASA Science}}\\[2pt]
{\tiny\url{https://science.nasa.gov/mission/europa-clipper/}}\\[4pt]
{\footnotesize Europa Clipper - NASA Science
Scientists predict Europa has a salty ocean beneath its icy crust that could hold the building blocks necessary to sustain life.TypeOrbiterLaunchOct. 10, 2024TargetEuropaObjectiveDetermine if Europa has conditions suitable to support lifeMission OverviewSpacecraftWhy Europa?What will Europa Clipper do?Europa Clipper’s main science goal is to determine whether there are places below the surface of Jupiter’s icy moon, Europa, that could support life.The mission’s three main science objectives are to understand the nature of the ice shell and the ocean beneath it, along with the moon’s composition and geology. The mission’s detailed exploration of Europa will help scientists better understand the astrobiological potential for habitable worlds beyond our planet.Read MoreAn artist's concept of NASA’s Europa Clipper spacecraft.NASA/JPL-CaltechWhat is NASA’s Europa Clipper spacecraft?Europa Clipper is a robotic solar-powered spacecraft built to conduct the first detailed investigations of Jupiter's icy moon Europa.With its solar arrays deployed, Europa Clipper spans more than 100 feet (about 30 meters) – about the length of a basketball court. The main body of the spacecraft consists of its avionics vault, radiofrequency module, and propulsion module.Read MoreWorkers use a crane to lift Europa Clipper's high-gain antenna as they prepare to install it on the spacecraft on Aug. 14, 2023. NASA/JPL-CaltechWhy Europa?The search for life beyond Earth is one of NASA’s primary objectives. If humans are to truly understand our place in the universe, we must learn whether our planet is the only place where life exists. So the search is on! There is strong evidence Jupiter's moon Europa has a saltwater ocean that may be one of the best places to look for environments where life could exist beyond Earth.Read MoreThis view of Jupiter’s icy moon Europa was captured by the JunoCam imager aboard NASA’s Juno spacecraft during the mission’s close flyby on Sept. 29, 2022.}
\end{docenv}
\begin{docenv}
{\footnotesize\textcolor{doccolor}{\textbf{Doc~3}}}\hfill{\tiny\texttt{score:~44.210}}\ {\tiny\textcolor{gray}{\{'relevant': \{'53573702-4140-4812-8658-ae10a54edfc8': \{'value': 'no', 'timestamp': 1723034566\}\}\}}}\\[3pt]
{\small\textit{NASA’s Europa Clipper Survives and Thrives in ‘Outer Space’ on Earth - NASA}}\\[2pt]
{\tiny\url{https://www.nasa.gov/missions/europa-clipper/nasas-europa-clipper-survives-and-thrives-in-outer-space-on-earth/}}\\[4pt]
{\footnotesize NASA’s Europa Clipper Survives and Thrives in ‘Outer Space’ on Earth - NASA
Europa Clipper’s main science goal is to determine whether there are places below the surface of Jupiter’s icy moon, Europa, that could support life. The mission’s three main science objectives are to determine the thickness of the moon’s icy shell and its surface interactions with the ocean below, to investigate its composition, and to characterize its geology. The mission’s detailed exploration of Europa will help scientists better understand the astrobiological potential for habitable worlds beyond our planet.
Managed by Caltech in Pasadena, California, JPL leads the development of the Europa Clipper mission in partnership with the Johns Hopkins Applied Physics Laboratory (APL) in Laurel, Maryland, for NASA’s Science Mission Directorate in Washington. APL designed the main spacecraft body in collaboration with JPL and NASA’s Goddard Space Flight Center in Greenbelt, Maryland. The Planetary Missions Program Office at NASA’s Marshall Space Flight Center in Huntsville, Alabama, executes program management of the Europa Clipper mission.
Find more information about Europa here:
europa.nasa.gov
News Media Contacts
Gretchen McCartneyJet Propulsion Laboratory, Pasadena, Calif.818-393-6215gretchen.p.mccartney@jpl.nasa.gov
Karen Fox / Charles BlueNASA Headquarters, Washington301-286-6284 / 202-802-5345karen.c.fox@nasa.gov / charles.e.blue@nasa.gov
2024-032

Share

Details

Last Updated

Mar 27, 2024

Related TermsEuropa ClipperEuropaJet Propulsion LaboratoryJupiterThe Solar System

 

Explore More

 

6 min read
New Images From Euclid Mission Reveal Wide View of the Dark Universe

Article

	2 days ago	

 

4 min read
NASA’s Psyche Fires Up Its Sci-Fi-Worthy Thrusters

Article

	3 days ago	

 

7 min read
NASA Tool Gets Ready to Image Faraway Planets

Article

	4 days ago	

 

Keep Exploring
Discover Related Topics

Missions

 

Humans in Space

 

Climate Change

 

Solar System

 

The National Aeronautics and Space Administration}
\end{docenv}
\begin{docenv}
{\footnotesize\textcolor{doccolor}{\textbf{Doc~4}}}\hfill{\tiny\texttt{score:~21.183}}\ {\tiny\textcolor{gray}{\{'relevant': \{'53573702-4140-4812-8658-ae10a54edfc8': \{'value': 'yes', 'timestamp': 1723034786\}\}\}}}\\[3pt]
{\small\textit{Surface of Jupiter's Moon Europa Churned by Small Impacts – NASA's Europa Clipper}}\\[2pt]
{\tiny\url{https://europa.nasa.gov/news/34/surface-of-jupiters-moon-europa-churned-by-small-impacts/}}\\[4pt]
{\footnotesize Surface of Jupiter's Moon Europa Churned by Small Impacts – NASA's Europa Clipper
Jupiter’s icy moon Europa may be the most promising place in the solar system to find present-day environments suitable for life beyond Earth. 

Europa: A World of Ice, With Potential for Life

Jupiter’s moon Europa has an icy crust covering a vast, global ocean. The rocky layer underneath may be hot enough to melt, leading to undersea volcanoes.

Europa's Interior May Be Hot Enough to Fuel Seafloor Volcanoes

Scientists have theorized on the origin of the water plumes possibly erupting from Jupiter's moon Europa. Recent research adds a new potential source to the mix.

Potential Plumes on Europa Could Come From Water in the Crust

New lab experiments re-create the environment of Europa and find that the icy moon shines, even on its nightside. The effect is more than just a cool visual.

Europa Glows: Radiation Does a Bright Number on Jupiter's Moon

The surface of Jupiter's moon Europa features a widely varied landscape, including ridges, bands, small rounded domes and disrupted spaces that geologists call "chaos terrain." Three newly reprocessed images, taken by NASA's Galileo spacecraft in the late 1990s, reveal details in diverse surface features on Europa.

Newly Reprocessed Images of Europa Show 'Chaos Terrain' in Crisp Detail

Scientists discovered that the yellow color visible on portions of the surface of Europa is actually sodium chloride, a compound known on Earth as table salt, which is also the principal component of sea salt.

Table Salt Compound Spotted on Europa

A full-scale prototype antenna, which at 10 feet tall is the same height as a standard basketball hoop, is in the Experimental Test Range (ETR) at NASA's Langley Research Center in Hampton, Virginia.

Europa Clipper High Gain Antenna Undergoes Testing at Langley

Students and families can meet a Mars rover, take a virtual tour through our solar system and explore alien worlds with NASA's Jet Propulsion Laboratory at the Clippers SciFest SoCal.

Catch NASA's JPL at the Clippers SciFest This Weekend

NASA has decided to replace the current magnetometer on the upcoming Europa Clipper mission with a less complex instrument. 

NASA Seeks New Options for Science Instrument on Europa Clipper}
\end{docenv}
\begin{docenv}
{\footnotesize\textcolor{doccolor}{\textbf{Doc~5}}}\hfill{\tiny\texttt{score:~20.931}}\ {\tiny\textcolor{gray}{\{'relevant': \{'53573702-4140-4812-8658-ae10a54edfc8': \{'value': 'yes', 'timestamp': 1723035044\}\}\}}}\\[3pt]
{\small\textit{Surface of Jupiter's Moon Europa Churned by Small Impacts – NASA's Europa Clipper}}\\[2pt]
{\tiny\url{https://europa.nasa.gov/news/34/surface-of-jupiters-moon-europa-churned-by-small-impacts/}}\\[4pt]
{\footnotesize Surface of Jupiter's Moon Europa Churned by Small Impacts – NASA's Europa Clipper
Surface of Jupiter's Moon Europa Churned by Small Impacts – NASA's Europa Clipper

NASA

Europa Clipper

Skip Navigation

menu

close modal

News

| July 12, 2021

Surface of Jupiter's Moon Europa Churned by Small Impacts

In this zoomed-in image of Europa’s surface, captured by NASA’s Galileo mission, the thin, bright layer, visible atop a cliff in the center shows the kind of areas churned by impact gardening. Credit: NASA/JPL-Caltech

Jupiter’s moon Europa and its global ocean may currently have conditions suitable for life. Scientists are studying processes on the icy surface as they prepare to explore.
It’s easy to see the impact of space debris on our Moon, where the ancient, battered surface is covered with craters and scars. Jupiter’s icy moon Europa withstands a similar trouncing – along with a punch of super-intense radiation. As the uppermost surface of the icy moon churns, material brought to the surface is zapped by high-energy electron radiation accelerated by Jupiter.
NASA-funded scientists are studying the cumulative effects of small impacts on Europa’s surface as they prepare to explore the distant moon with the Europa Clipper mission and study the possibilities for a future lander mission. Europa is of particular scientific interest because its salty ocean, which lies beneath a thick layer of ice, may currently have conditions suitable for existing life. That water may even make its way into the icy crust and onto the moon’s surface.
New research and modeling estimate how far down that surface is disturbed by the process called “impact gardening.” The work, published July 12 in Nature Astronomy, estimates that the surface of Europa has been churned by small impacts to an average depth of about 12 inches (30 centimeters) over tens of millions of years. And any molecules that might qualify as potential biosignatures, which include chemical signs of life, could be affected at that depth.
That’s because the impacts would churn some material to the surface, where radiation would likely break the bonds of any potential large, delicate molecules generated by biology.}
\end{docenv}
\begin{docenv}
{\footnotesize\textcolor{doccolor}{\textbf{Doc~6}}}\hfill{\tiny\texttt{score:~20.844}}\ {\tiny\textcolor{gray}{\{'relevant': \{'53573702-4140-4812-8658-ae10a54edfc8': \{'value': 'yes', 'timestamp': 1723035037\}\}\}}}\\[3pt]
{\small\textit{Europa's Interior May Be Hot Enough to Fuel Seafloor Volcanoes – NASA's Europa Clipper}}\\[2pt]
{\tiny\url{https://europa.nasa.gov/news/32/europas-interior-may-be-hot-enough-to-fuel-seafloor-volcanoes/}}\\[4pt]
{\footnotesize Europa's Interior May Be Hot Enough to Fuel Seafloor Volcanoes – NASA's Europa Clipper
Greenland Ice, Jupiter Moon Share Similar Feature

Jupiter’s moon Europa and its global ocean may currently have conditions suitable for life. Scientists are studying processes on the icy surface as they prepare to explore.

Surface of Jupiter's Moon Europa Churned by Small Impacts

Jupiter’s icy moon Europa may be the most promising place in the solar system to find present-day environments suitable for life beyond Earth. 

Europa: A World of Ice, With Potential for Life

Scientists have theorized on the origin of the water plumes possibly erupting from Jupiter's moon Europa. Recent research adds a new potential source to the mix.

Potential Plumes on Europa Could Come From Water in the Crust

New lab experiments re-create the environment of Europa and find that the icy moon shines, even on its nightside. The effect is more than just a cool visual.

Europa Glows: Radiation Does a Bright Number on Jupiter's Moon

The surface of Jupiter's moon Europa features a widely varied landscape, including ridges, bands, small rounded domes and disrupted spaces that geologists call "chaos terrain." Three newly reprocessed images, taken by NASA's Galileo spacecraft in the late 1990s, reveal details in diverse surface features on Europa.

Newly Reprocessed Images of Europa Show 'Chaos Terrain' in Crisp Detail

Scientists discovered that the yellow color visible on portions of the surface of Europa is actually sodium chloride, a compound known on Earth as table salt, which is also the principal component of sea salt.

Table Salt Compound Spotted on Europa

A full-scale prototype antenna, which at 10 feet tall is the same height as a standard basketball hoop, is in the Experimental Test Range (ETR) at NASA's Langley Research Center in Hampton, Virginia.

Europa Clipper High Gain Antenna Undergoes Testing at Langley

Students and families can meet a Mars rover, take a virtual tour through our solar system and explore alien worlds with NASA's Jet Propulsion Laboratory at the Clippers SciFest SoCal.

Catch NASA's JPL at the Clippers SciFest This Weekend

NASA has decided to replace the current magnetometer on the upcoming Europa Clipper mission with a less complex instrument.}
\end{docenv}
\begin{refanswerenv}
{\footnotesize\textit{\textcolor{gray}{\textbf{Reference answer:}\enspace The surface of Jupiter's moon, Europa, is made up of ice. It is covered in a thick crust of frozen water that is estimated to be between 10 to 100 kilometers thick. The surface also shows signs of past and present geological activity, such as ridges, chaos terrain, and lenticulae. These features suggest that there is heat being generated from within Europa, which is thought to be due to tidal heating caused by Jupiter's gravitational pull. Additionally, the surface of Europa is subject to bombardment by radiation from Jupiter, which can cause the ice to darken and form a reddish-brown color in some areas.}}}
\end{refanswerenv}
\vspace{0.5em}
\noindent{\small\textcolor{usercolor}{\textbf{$\triangleright$\enspace User --- Turn~4}}}\hfill{\tiny\textcolor{gray}{Factoid $\cdot$ Follow-up $\cdot$ ANSWERABLE}}
\begin{userbubble}
Who funds the Clipper mission?
\end{userbubble}
\noindent\hfill{\small\textcolor{agentcolor}{\textbf{Agent\enspace$\triangleleft$}}}
\begin{agentbubble}
The Europa Clipper mission is sponsored by NASA. The total contract award amount for launch services is approximately \$178 million.
\end{agentbubble}
\noindent\hfill{\footnotesize\textcolor{doccolor}{\textbf{$\hookrightarrow$\enspace Retrieved documents~(4)}}}
\begin{docenv}
{\footnotesize\textcolor{doccolor}{\textbf{Doc~1}}}\hfill{\tiny\texttt{score:~48.076}}\ {\tiny\textcolor{gray}{\{'relevant': \{'53573702-4140-4812-8658-ae10a54edfc8': \{'value': 'no', 'timestamp': 1723035700\}\}\}}}\\[3pt]
{\small\textit{Science | Mission – NASA's Europa Clipper}}\\[2pt]
{\tiny\url{https://europa.nasa.gov/mission/science/}}\\[4pt]
{\footnotesize Science | Mission – NASA's Europa Clipper
Science | Mission – NASA's Europa Clipper

NASA

Europa Clipper

Skip Navigation

menu

close modal

Europa Clipper Mission

Science

Mission

MENU

Overview

Science

Timeline

The Team

Overview

History

FAQ

Key Fact

What is Europa Clipper's main science goal? 

Europa Clipper’s main science goal is to determine whether there are places below the surface of Jupiter’s icy moon, Europa, that could support life.

OverviewOverview

Deep beneath the crust of Jupiter’s frozen moon Europa lies a massive liquid water ocean. Exploring this ocean world with our Europa Clipper spacecraft could provide new clues in our search for life beyond Earth. Scheduled to launch in October 2024, Europa Clipper’s main science goal is to determine whether there are places below the surface of Jupiter’s icy moon, Europa, that could support life. Credit: NASA

Jupiter’s moon Europa shows strong evidence for an ocean of liquid water beneath its icy crust. Beyond Earth, Europa is considered one of the most promising currently habitable environments in our solar system. 
Below Europa’s icy surface, evidence suggests there is a global ocean with more water than all of Earth’s oceans combined. Europa could have all the “ingredients” needed for life as we know it: 

Water: Twice as much as Earth’s oceans 

Organics: Essential chemical building blocks from a variety of sources

Energy: Chemical energy sources from the surface and the sea floor

Stability: Conditions remaining similar for 4 billion years

NASA's Europa Clipper spacecraft will perform approximately 50 close flybys of the moon, gathering detailed measurements to investigate whether the moon could have conditions suitable for life. Europa Clipper is not a life detection mission – its main science goal is to determine whether there are places below Europa’s surface that could support life.

Europa Up Close

Take an in-depth look at Jupiter's icy moon and Europa Clipper's plans for investigation.

Learn more ›

Science ObjectivesScience Objectives
Europa Clipper has three main science objectives: 

Determine the Thickness of Europa’s Icy Shell \& How the Ocean Interacts With the Surface}
\end{docenv}
\begin{docenv}
{\footnotesize\textcolor{doccolor}{\textbf{Doc~2}}}\hfill{\tiny\texttt{score:~46.229}}\ {\tiny\textcolor{gray}{\{'relevant': \{'53573702-4140-4812-8658-ae10a54edfc8': \{'value': 'no', 'timestamp': 1723035735\}\}\}}}\\[3pt]
{\small\textit{Europa Clipper - NASA Science}}\\[2pt]
{\tiny\url{https://science.nasa.gov/mission/europa-clipper/}}\\[4pt]
{\footnotesize Europa Clipper - NASA Science
Scientists predict Europa has a salty ocean beneath its icy crust that could hold the building blocks necessary to sustain life.TypeOrbiterLaunchOct. 10, 2024TargetEuropaObjectiveDetermine if Europa has conditions suitable to support lifeMission OverviewSpacecraftWhy Europa?What will Europa Clipper do?Europa Clipper’s main science goal is to determine whether there are places below the surface of Jupiter’s icy moon, Europa, that could support life.The mission’s three main science objectives are to understand the nature of the ice shell and the ocean beneath it, along with the moon’s composition and geology. The mission’s detailed exploration of Europa will help scientists better understand the astrobiological potential for habitable worlds beyond our planet.Read MoreAn artist's concept of NASA’s Europa Clipper spacecraft.NASA/JPL-CaltechWhat is NASA’s Europa Clipper spacecraft?Europa Clipper is a robotic solar-powered spacecraft built to conduct the first detailed investigations of Jupiter's icy moon Europa.With its solar arrays deployed, Europa Clipper spans more than 100 feet (about 30 meters) – about the length of a basketball court. The main body of the spacecraft consists of its avionics vault, radiofrequency module, and propulsion module.Read MoreWorkers use a crane to lift Europa Clipper's high-gain antenna as they prepare to install it on the spacecraft on Aug. 14, 2023. NASA/JPL-CaltechWhy Europa?The search for life beyond Earth is one of NASA’s primary objectives. If humans are to truly understand our place in the universe, we must learn whether our planet is the only place where life exists. So the search is on! There is strong evidence Jupiter's moon Europa has a saltwater ocean that may be one of the best places to look for environments where life could exist beyond Earth.Read MoreThis view of Jupiter’s icy moon Europa was captured by the JunoCam imager aboard NASA’s Juno spacecraft during the mission’s close flyby on Sept. 29, 2022.}
\end{docenv}
\begin{docenv}
{\footnotesize\textcolor{doccolor}{\textbf{Doc~3}}}\hfill{\tiny\texttt{score:~46.110}}\ {\tiny\textcolor{gray}{\{'relevant': \{'53573702-4140-4812-8658-ae10a54edfc8': \{'value': 'no', 'timestamp': 1723035776\}\}\}}}\\[3pt]
{\small\textit{Europa Clipper High Gain Antenna Undergoes Testing at Langley – NASA's Europa Clipper}}\\[2pt]
{\tiny\url{https://europa.nasa.gov/news/23/europa-clipper-high-gain-antenna-undergoes-testing-at-langley/}}\\[4pt]
{\footnotesize Europa Clipper High Gain Antenna Undergoes Testing at Langley – NASA's Europa Clipper
Jupiter’s moon Europa and its global ocean may currently have conditions suitable for life. Scientists are studying processes on the icy surface as they prepare to explore.

Surface of Jupiter's Moon Europa Churned by Small Impacts

Jupiter’s icy moon Europa may be the most promising place in the solar system to find present-day environments suitable for life beyond Earth. 

Europa: A World of Ice, With Potential for Life

Jupiter’s moon Europa has an icy crust covering a vast, global ocean. The rocky layer underneath may be hot enough to melt, leading to undersea volcanoes.

Europa's Interior May Be Hot Enough to Fuel Seafloor Volcanoes

Jupiter’s moon Europa may have the potential to harbor life. The spacecraft will use multiple flybys of the moon to investigate the habitability of this ocean world.

NASA's Europa Clipper Builds Hardware, Moves Toward Assembly

In an unprecedented year, the Europa Clipper project continues on the path toward unprecedented science. 

On the Path Towards Unprecedented Science

Scientists have theorized on the origin of the water plumes possibly erupting from Jupiter's moon Europa. Recent research adds a new potential source to the mix.

Potential Plumes on Europa Could Come From Water in the Crust

New lab experiments re-create the environment of Europa and find that the icy moon shines, even on its nightside. The effect is more than just a cool visual.

Europa Glows: Radiation Does a Bright Number on Jupiter's Moon

The surface of Jupiter's moon Europa features a widely varied landscape, including ridges, bands, small rounded domes and disrupted spaces that geologists call "chaos terrain." Three newly reprocessed images, taken by NASA's Galileo spacecraft in the late 1990s, reveal details in diverse surface features on Europa.

Newly Reprocessed Images of Europa Show 'Chaos Terrain' in Crisp Detail

An icy ocean world in our solar system that could tell us more about the potential for life on other worlds is coming into focus with confirmation of the Europa Clipper mission's next phase. 

Mission to Jupiter's Icy Moon Confirmed}
\end{docenv}
\begin{docenv}
{\footnotesize\textcolor{doccolor}{\textbf{Doc~4}}}\hfill{\tiny\texttt{score:~16.573}}\ {\tiny\textcolor{gray}{\{'relevant': \{'53573702-4140-4812-8658-ae10a54edfc8': \{'value': 'yes', 'timestamp': 1723036431\}\}\}}}\\[3pt]
{\small\textit{NASA Awards Launch Services Contract for Europa Clipper Mission – NASA's Europa Clipper}}\\[2pt]
{\tiny\url{https://europa.nasa.gov/news/36/nasa-awards-launch-services-contract-for-europa-clipper-mission/}}\\[4pt]
{\footnotesize NASA Awards Launch Services Contract for Europa Clipper Mission – NASA's Europa Clipper
NASA Awards Launch Services Contract for Europa Clipper Mission – NASA's Europa Clipper

NASA

Europa Clipper

Skip Navigation

menu

close modal

News

| July 23, 2021

NASA Awards Launch Services Contract for Europa Clipper Mission

This illustration, updated in December 2020, depicts NASA's Europa Clipper spacecraft. With an internal global ocean twice the size of Earth's oceans combined, Jupiter's moon Europa may have the potential to harbor life. The Europa Clipper orbiter will swoop around Jupiter on an elliptical path, dipping close to the moon on each flyby to collect data. Credits: NASA/JPL-Caltech

NASA has selected Space Exploration Technologies Corp. (SpaceX) of Hawthorne, California, to provide launch services for Earth’s first mission to conduct detailed investigations of Jupiter's moon Europa.
The Europa Clipper mission will launch in October 2024 on a Falcon Heavy rocket from Launch Complex 39A at NASA’s Kennedy Space Center in Florida. The total contract award amount for launch services is approximately \$178 million.
Europa Clipper will conduct a detailed survey of Europa and use a sophisticated suite of science instruments to investigate whether the icy moon has conditions suitable for life. Key mission objectives are to produce high-resolution images of Europa's surface, determine its composition, look for signs of recent or ongoing geological activity, measure the thickness of the moon’s icy shell, search for subsurface lakes, and determine the depth and salinity of Europa's ocean.
NASA’s Launch Services Program at Kennedy will manage the Europa Clipper launch service. NASA's Jet Propulsion Laboratory in Southern California leads the development of the Europa Clipper mission in partnership with the Johns Hopkins Applied Physics Laboratory in Laurel, Maryland, for NASA’s Science Mission Directorate in Washington. The Planetary Missions Program Office at NASA’s Marshall Space Flight Center in Huntsville, Alabama, executes program management of the Europa Clipper mission.
For more information about the Europa Clipper mission, visit:
https://www.nasa.gov/europa

News Media Contacts
Monica Witt
Headquarters, Washington
202-358-1100
monica.j.witt@nasa.gov}
\end{docenv}
\begin{refanswerenv}
{\footnotesize\textit{\textcolor{gray}{\textbf{Reference answer:}\enspace The Europa Clipper mission is funded by NASA, which is a U.S. government agency responsible for civilian space exploration and aeronautics research. NASA's budget comes primarily from Congressional appropriations.}}}
\end{refanswerenv}
\vspace{0.5em}
\noindent{\small\textcolor{usercolor}{\textbf{$\triangleright$\enspace User --- Turn~5}}}\hfill{\tiny\textcolor{gray}{Factoid, Summarization $\cdot$ Follow-up $\cdot$ ANSWERABLE}}
\begin{userbubble}
Does NASA have an intern program?
\end{userbubble}
\noindent\hfill{\small\textcolor{agentcolor}{\textbf{Agent\enspace$\triangleleft$}}}
\begin{agentbubble}
Yes, NASA does offer an internship program. Interns may support a variety of business and program management work, including human resources, administration and program analysis, accounting and budgeting, communication and public relations, procurement and contracts, information technology, security, and education.
\end{agentbubble}
\noindent\hfill{\footnotesize\textcolor{doccolor}{\textbf{$\hookrightarrow$\enspace Retrieved documents~(5)}}}
\begin{docenv}
{\footnotesize\textcolor{doccolor}{\textbf{Doc~1}}}\hfill{\tiny\texttt{score:~47.200}}\ {\tiny\textcolor{gray}{\{'relevant': \{'53573702-4140-4812-8658-ae10a54edfc8': \{'value': 'no', 'timestamp': 1723036835\}\}\}}}\\[3pt]
{\small\textit{Science | Mission – NASA's Europa Clipper}}\\[2pt]
{\tiny\url{https://europa.nasa.gov/mission/science/}}\\[4pt]
{\footnotesize Science | Mission – NASA's Europa Clipper
Science | Mission – NASA's Europa Clipper

NASA

Europa Clipper

Skip Navigation

menu

close modal

Europa Clipper Mission

Science

Mission

MENU

Overview

Science

Timeline

The Team

Overview

History

FAQ

Key Fact

What is Europa Clipper's main science goal? 

Europa Clipper’s main science goal is to determine whether there are places below the surface of Jupiter’s icy moon, Europa, that could support life.

OverviewOverview

Deep beneath the crust of Jupiter’s frozen moon Europa lies a massive liquid water ocean. Exploring this ocean world with our Europa Clipper spacecraft could provide new clues in our search for life beyond Earth. Scheduled to launch in October 2024, Europa Clipper’s main science goal is to determine whether there are places below the surface of Jupiter’s icy moon, Europa, that could support life. Credit: NASA

Jupiter’s moon Europa shows strong evidence for an ocean of liquid water beneath its icy crust. Beyond Earth, Europa is considered one of the most promising currently habitable environments in our solar system. 
Below Europa’s icy surface, evidence suggests there is a global ocean with more water than all of Earth’s oceans combined. Europa could have all the “ingredients” needed for life as we know it: 

Water: Twice as much as Earth’s oceans 

Organics: Essential chemical building blocks from a variety of sources

Energy: Chemical energy sources from the surface and the sea floor

Stability: Conditions remaining similar for 4 billion years

NASA's Europa Clipper spacecraft will perform approximately 50 close flybys of the moon, gathering detailed measurements to investigate whether the moon could have conditions suitable for life. Europa Clipper is not a life detection mission – its main science goal is to determine whether there are places below Europa’s surface that could support life.

Europa Up Close

Take an in-depth look at Jupiter's icy moon and Europa Clipper's plans for investigation.

Learn more ›

Science ObjectivesScience Objectives
Europa Clipper has three main science objectives: 

Determine the Thickness of Europa’s Icy Shell \& How the Ocean Interacts With the Surface}
\end{docenv}
\begin{docenv}
{\footnotesize\textcolor{doccolor}{\textbf{Doc~2}}}\hfill{\tiny\texttt{score:~45.433}}\ {\tiny\textcolor{gray}{\{'relevant': \{'53573702-4140-4812-8658-ae10a54edfc8': \{'value': 'no', 'timestamp': 1723036852\}\}\}}}\\[3pt]
{\small\textit{Europa Clipper - NASA Science}}\\[2pt]
{\tiny\url{https://science.nasa.gov/mission/europa-clipper/}}\\[4pt]
{\footnotesize Europa Clipper - NASA Science
Scientists predict Europa has a salty ocean beneath its icy crust that could hold the building blocks necessary to sustain life.TypeOrbiterLaunchOct. 10, 2024TargetEuropaObjectiveDetermine if Europa has conditions suitable to support lifeMission OverviewSpacecraftWhy Europa?What will Europa Clipper do?Europa Clipper’s main science goal is to determine whether there are places below the surface of Jupiter’s icy moon, Europa, that could support life.The mission’s three main science objectives are to understand the nature of the ice shell and the ocean beneath it, along with the moon’s composition and geology. The mission’s detailed exploration of Europa will help scientists better understand the astrobiological potential for habitable worlds beyond our planet.Read MoreAn artist's concept of NASA’s Europa Clipper spacecraft.NASA/JPL-CaltechWhat is NASA’s Europa Clipper spacecraft?Europa Clipper is a robotic solar-powered spacecraft built to conduct the first detailed investigations of Jupiter's icy moon Europa.With its solar arrays deployed, Europa Clipper spans more than 100 feet (about 30 meters) – about the length of a basketball court. The main body of the spacecraft consists of its avionics vault, radiofrequency module, and propulsion module.Read MoreWorkers use a crane to lift Europa Clipper's high-gain antenna as they prepare to install it on the spacecraft on Aug. 14, 2023. NASA/JPL-CaltechWhy Europa?The search for life beyond Earth is one of NASA’s primary objectives. If humans are to truly understand our place in the universe, we must learn whether our planet is the only place where life exists. So the search is on! There is strong evidence Jupiter's moon Europa has a saltwater ocean that may be one of the best places to look for environments where life could exist beyond Earth.Read MoreThis view of Jupiter’s icy moon Europa was captured by the JunoCam imager aboard NASA’s Juno spacecraft during the mission’s close flyby on Sept. 29, 2022.}
\end{docenv}
\begin{docenv}
{\footnotesize\textcolor{doccolor}{\textbf{Doc~3}}}\hfill{\tiny\texttt{score:~44.823}}\ {\tiny\textcolor{gray}{\{'relevant': \{'53573702-4140-4812-8658-ae10a54edfc8': \{'value': 'no', 'timestamp': 1723036923\}\}\}}}\\[3pt]
{\small\textit{Europa Clipper High Gain Antenna Undergoes Testing at Langley – NASA's Europa Clipper}}\\[2pt]
{\tiny\url{https://europa.nasa.gov/news/23/europa-clipper-high-gain-antenna-undergoes-testing-at-langley/}}\\[4pt]
{\footnotesize Europa Clipper High Gain Antenna Undergoes Testing at Langley – NASA's Europa Clipper
Jupiter’s moon Europa and its global ocean may currently have conditions suitable for life. Scientists are studying processes on the icy surface as they prepare to explore.

Surface of Jupiter's Moon Europa Churned by Small Impacts

Jupiter’s icy moon Europa may be the most promising place in the solar system to find present-day environments suitable for life beyond Earth. 

Europa: A World of Ice, With Potential for Life

Jupiter’s moon Europa has an icy crust covering a vast, global ocean. The rocky layer underneath may be hot enough to melt, leading to undersea volcanoes.

Europa's Interior May Be Hot Enough to Fuel Seafloor Volcanoes

Jupiter’s moon Europa may have the potential to harbor life. The spacecraft will use multiple flybys of the moon to investigate the habitability of this ocean world.

NASA's Europa Clipper Builds Hardware, Moves Toward Assembly

In an unprecedented year, the Europa Clipper project continues on the path toward unprecedented science. 

On the Path Towards Unprecedented Science

Scientists have theorized on the origin of the water plumes possibly erupting from Jupiter's moon Europa. Recent research adds a new potential source to the mix.

Potential Plumes on Europa Could Come From Water in the Crust

New lab experiments re-create the environment of Europa and find that the icy moon shines, even on its nightside. The effect is more than just a cool visual.

Europa Glows: Radiation Does a Bright Number on Jupiter's Moon

The surface of Jupiter's moon Europa features a widely varied landscape, including ridges, bands, small rounded domes and disrupted spaces that geologists call "chaos terrain." Three newly reprocessed images, taken by NASA's Galileo spacecraft in the late 1990s, reveal details in diverse surface features on Europa.

Newly Reprocessed Images of Europa Show 'Chaos Terrain' in Crisp Detail

An icy ocean world in our solar system that could tell us more about the potential for life on other worlds is coming into focus with confirmation of the Europa Clipper mission's next phase. 

Mission to Jupiter's Icy Moon Confirmed}
\end{docenv}
\begin{docenv}
{\footnotesize\textcolor{doccolor}{\textbf{Doc~4}}}\hfill{\tiny\texttt{score:~22.208}}\ {\tiny\textcolor{gray}{\{'relevant': \{'53573702-4140-4812-8658-ae10a54edfc8': \{'value': 'yes', 'timestamp': 1723037846\}, '9dcd142c-415e-4161-a82f-c8b14bfb36ac': \{'value': 'yes', 'timestamp': 1726570113\}\}\}}}\\[3pt]
{\small\textit{NASA Internship Programs - NASA}}\\[2pt]
{\tiny\url{https://intern.nasa.gov/}}\\[4pt]
{\footnotesize NASA Internship Programs - NASA
Launch your career with a Pathways internship.
NASA’s work-study (co-op) program is a starting point to a career at NASA. Pathways interns gain valuable work experience and professional development. Paid federal civil servant opportunities with benefits are offered across most NASA facilities. Completion of the Pathways program may lead to a NASA job offer.

	Learn More 	

 

Fellowship

NASA Fellowships allow graduate-level students to pursue research projects in response to the agency’s current research priorities.
NASA Fellowships support graduate-level projects and proposals which contribute to a NASA research opportunity. Student researchers are exposed to NASA’s innovation-oriented culture and facilities. Fellows participate in prestigious conferences and a center based research experience. Fellowships support academic institutions by enhancing graduate learning and development.

	Learn More 	

 

Internships at NASA’s Jet Propulsion Laboratory

Discover exciting internships and research opportunities at the leading center for robotic exploration of the solar system.
An internship at NASA’s Jet Propulsion Laboratory is a chance to do the impossible. Our internships put you right in the action with the scientists and engineers who’ve helped make JPL the leading center for robotic exploration of the solar system. Our programs are as varied as the places we explore, with opportunities across the STEM spectrum for undergrads, graduate students, postdocs and faculty. Join us and do something out of this world.

	Learn More 	

 

International Internship

NASA Intern and Fellow Opportunities for International Students
NASA International Internship (NASA I²) seeks to better prepare students to work in a global environment and on multicultural international missions. NASA and the nation benefit from a cadre of future scientists, engineers and other professionals who become familiar and experienced in multinational environments. Internship sessions are arranged in three sessions during the calendar year (spring, summer and fall). 

	Learn More 	

 

Eligibility at a Glance}
\end{docenv}
\begin{docenv}
{\footnotesize\textcolor{doccolor}{\textbf{Doc~5}}}\hfill{\tiny\texttt{score:~22.040}}\ {\tiny\textcolor{gray}{\{'relevant': \{'53573702-4140-4812-8658-ae10a54edfc8': \{'value': 'yes', 'timestamp': 1723037921\}\}\}}}\\[3pt]
{\small\textit{NASA Internship Programs - NASA}}\\[2pt]
{\tiny\url{https://intern.nasa.gov/}}\\[4pt]
{\footnotesize NASA Internship Programs - NASA
At NASA, our engineers, mathematicians, scientists, accountants, writers, IT specialists, project managers, program analysts, and many other professionals work together to break barriers to achieve the seemingly impossible. Non-engineering interns may support a variety of business and program management work, including procurement, budgeting, accounting, information technology, and security.

Frequently Asked Questions

 

Frequently Asked Questions

	More Frequently Asked Questions 	

 

Are Internships Paid?

The majority of interns receive a stipend award, but there are some volunteer opportunities noted in project descriptions.

 

Is housing my responsibility if I receive an internship?

Interns are responsible for making their own housing arrangements. This includes locating their own housing options and paying for their housing. Centers may be able to offer minimal assistance by providing a list of local available housing and/or establishing a private social media group for interns to utilize for relocation planning purposes.

 

Is a Letter of Recommendation required?

The Office of STEM Engagement does not require a letter of recommendation. 

 

Will I need a car?

Your center location and housing options will determine if you need a car. At most centers, unless you find housing within walking distance, you will most likely need a car. However, there are ample public transportation options at some centers

 

I'm not a science or engineering student. Can I still get an internship at NASA?

Yes! Interns may support a variety of business and program management work, including human resources, administration and program analysis, accounting and budget, communication and public relations, procurement and contracts, information technology, security, and education. No matter your academic major, you’ll find that NASA is a place where you’ll belong to a supportive community that values your background, perspective, and life experiences.

 

My Path to NASA

Astronauts are often the face of NASA, but our workforce is made up of all kinds of professionals. We are scientists, engineers, IT specialists, human resources specialists, accountants, writers, technicians and many other kinds of people working together. No matter what your strengths are, there's a place at NASA for you!

	Explore Careers at NASA}
\end{docenv}
\begin{refanswerenv}
{\footnotesize\textit{\textcolor{gray}{\textbf{Reference answer:}\enspace Yes, NASA does offer an internship program. The NASA Internships program is for high school, undergraduate, and graduate students, as well as educators, and provides students with the opportunity to work with professionals in their field of interest and gain valuable experience. You can find more information about the program and how to apply on NASA's official website.}}}
\end{refanswerenv}
\vspace{0.5em}
\noindent{\small\textcolor{usercolor}{\textbf{$\triangleright$\enspace User --- Turn~6}}}\hfill{\tiny\textcolor{gray}{Factoid, Summarization $\cdot$ Follow-up $\cdot$ ANSWERABLE}}
\begin{userbubble}
Does Clipper have an antennae?
\end{userbubble}
\noindent\hfill{\small\textcolor{agentcolor}{\textbf{Agent\enspace$\triangleleft$}}}
\begin{agentbubble}
Yes, the Europa Clipper spacecraft has an high-gain antenna. It is used for communication with Earth and sending data back to NASA. Workers were seen installing it on the spacecraft on August 14, 2023. Europa Clipper’s 10-foot-diameter (3-meter-diameter) high-gain antenna (HGA) and other smaller radio antennas will be used to receive commands from Earth through NASA’s Deep Space Network. 
\end{agentbubble}
\noindent\hfill{\footnotesize\textcolor{doccolor}{\textbf{$\hookrightarrow$\enspace Retrieved documents~(5)}}}
\begin{docenv}
{\footnotesize\textcolor{doccolor}{\textbf{Doc~1}}}\hfill{\tiny\texttt{score:~47.200}}\ {\tiny\textcolor{gray}{\{'relevant': \{'53573702-4140-4812-8658-ae10a54edfc8': \{'value': 'no', 'timestamp': 1723039478\}\}\}}}\\[3pt]
{\small\textit{Science | Mission – NASA's Europa Clipper}}\\[2pt]
{\tiny\url{https://europa.nasa.gov/mission/science/}}\\[4pt]
{\footnotesize Science | Mission – NASA's Europa Clipper
Science | Mission – NASA's Europa Clipper

NASA

Europa Clipper

Skip Navigation

menu

close modal

Europa Clipper Mission

Science

Mission

MENU

Overview

Science

Timeline

The Team

Overview

History

FAQ

Key Fact

What is Europa Clipper's main science goal? 

Europa Clipper’s main science goal is to determine whether there are places below the surface of Jupiter’s icy moon, Europa, that could support life.

OverviewOverview

Deep beneath the crust of Jupiter’s frozen moon Europa lies a massive liquid water ocean. Exploring this ocean world with our Europa Clipper spacecraft could provide new clues in our search for life beyond Earth. Scheduled to launch in October 2024, Europa Clipper’s main science goal is to determine whether there are places below the surface of Jupiter’s icy moon, Europa, that could support life. Credit: NASA

Jupiter’s moon Europa shows strong evidence for an ocean of liquid water beneath its icy crust. Beyond Earth, Europa is considered one of the most promising currently habitable environments in our solar system. 
Below Europa’s icy surface, evidence suggests there is a global ocean with more water than all of Earth’s oceans combined. Europa could have all the “ingredients” needed for life as we know it: 

Water: Twice as much as Earth’s oceans 

Organics: Essential chemical building blocks from a variety of sources

Energy: Chemical energy sources from the surface and the sea floor

Stability: Conditions remaining similar for 4 billion years

NASA's Europa Clipper spacecraft will perform approximately 50 close flybys of the moon, gathering detailed measurements to investigate whether the moon could have conditions suitable for life. Europa Clipper is not a life detection mission – its main science goal is to determine whether there are places below Europa’s surface that could support life.

Europa Up Close

Take an in-depth look at Jupiter's icy moon and Europa Clipper's plans for investigation.

Learn more ›

Science ObjectivesScience Objectives
Europa Clipper has three main science objectives: 

Determine the Thickness of Europa’s Icy Shell \& How the Ocean Interacts With the Surface}
\end{docenv}
\begin{docenv}
{\footnotesize\textcolor{doccolor}{\textbf{Doc~2}}}\hfill{\tiny\texttt{score:~45.433}}\ {\tiny\textcolor{gray}{\{'relevant': \{'53573702-4140-4812-8658-ae10a54edfc8': \{'value': 'yes', 'timestamp': 1723039506\}\}\}}}\\[3pt]
{\small\textit{Europa Clipper - NASA Science}}\\[2pt]
{\tiny\url{https://science.nasa.gov/mission/europa-clipper/}}\\[4pt]
{\footnotesize Europa Clipper - NASA Science
Scientists predict Europa has a salty ocean beneath its icy crust that could hold the building blocks necessary to sustain life.TypeOrbiterLaunchOct. 10, 2024TargetEuropaObjectiveDetermine if Europa has conditions suitable to support lifeMission OverviewSpacecraftWhy Europa?What will Europa Clipper do?Europa Clipper’s main science goal is to determine whether there are places below the surface of Jupiter’s icy moon, Europa, that could support life.The mission’s three main science objectives are to understand the nature of the ice shell and the ocean beneath it, along with the moon’s composition and geology. The mission’s detailed exploration of Europa will help scientists better understand the astrobiological potential for habitable worlds beyond our planet.Read MoreAn artist's concept of NASA’s Europa Clipper spacecraft.NASA/JPL-CaltechWhat is NASA’s Europa Clipper spacecraft?Europa Clipper is a robotic solar-powered spacecraft built to conduct the first detailed investigations of Jupiter's icy moon Europa.With its solar arrays deployed, Europa Clipper spans more than 100 feet (about 30 meters) – about the length of a basketball court. The main body of the spacecraft consists of its avionics vault, radiofrequency module, and propulsion module.Read MoreWorkers use a crane to lift Europa Clipper's high-gain antenna as they prepare to install it on the spacecraft on Aug. 14, 2023. NASA/JPL-CaltechWhy Europa?The search for life beyond Earth is one of NASA’s primary objectives. If humans are to truly understand our place in the universe, we must learn whether our planet is the only place where life exists. So the search is on! There is strong evidence Jupiter's moon Europa has a saltwater ocean that may be one of the best places to look for environments where life could exist beyond Earth.Read MoreThis view of Jupiter’s icy moon Europa was captured by the JunoCam imager aboard NASA’s Juno spacecraft during the mission’s close flyby on Sept. 29, 2022.}
\end{docenv}
\begin{docenv}
{\footnotesize\textcolor{doccolor}{\textbf{Doc~3}}}\hfill{\tiny\texttt{score:~44.823}}\ {\tiny\textcolor{gray}{\{'relevant': \{'53573702-4140-4812-8658-ae10a54edfc8': \{'value': 'no', 'timestamp': 1723039573\}\}\}}}\\[3pt]
{\small\textit{Europa Clipper High Gain Antenna Undergoes Testing at Langley – NASA's Europa Clipper}}\\[2pt]
{\tiny\url{https://europa.nasa.gov/news/23/europa-clipper-high-gain-antenna-undergoes-testing-at-langley/}}\\[4pt]
{\footnotesize Europa Clipper High Gain Antenna Undergoes Testing at Langley – NASA's Europa Clipper
Jupiter’s moon Europa and its global ocean may currently have conditions suitable for life. Scientists are studying processes on the icy surface as they prepare to explore.

Surface of Jupiter's Moon Europa Churned by Small Impacts

Jupiter’s icy moon Europa may be the most promising place in the solar system to find present-day environments suitable for life beyond Earth. 

Europa: A World of Ice, With Potential for Life

Jupiter’s moon Europa has an icy crust covering a vast, global ocean. The rocky layer underneath may be hot enough to melt, leading to undersea volcanoes.

Europa's Interior May Be Hot Enough to Fuel Seafloor Volcanoes

Jupiter’s moon Europa may have the potential to harbor life. The spacecraft will use multiple flybys of the moon to investigate the habitability of this ocean world.

NASA's Europa Clipper Builds Hardware, Moves Toward Assembly

In an unprecedented year, the Europa Clipper project continues on the path toward unprecedented science. 

On the Path Towards Unprecedented Science

Scientists have theorized on the origin of the water plumes possibly erupting from Jupiter's moon Europa. Recent research adds a new potential source to the mix.

Potential Plumes on Europa Could Come From Water in the Crust

New lab experiments re-create the environment of Europa and find that the icy moon shines, even on its nightside. The effect is more than just a cool visual.

Europa Glows: Radiation Does a Bright Number on Jupiter's Moon

The surface of Jupiter's moon Europa features a widely varied landscape, including ridges, bands, small rounded domes and disrupted spaces that geologists call "chaos terrain." Three newly reprocessed images, taken by NASA's Galileo spacecraft in the late 1990s, reveal details in diverse surface features on Europa.

Newly Reprocessed Images of Europa Show 'Chaos Terrain' in Crisp Detail

An icy ocean world in our solar system that could tell us more about the potential for life on other worlds is coming into focus with confirmation of the Europa Clipper mission's next phase. 

Mission to Jupiter's Icy Moon Confirmed}
\end{docenv}
\begin{docenv}
{\footnotesize\textcolor{doccolor}{\textbf{Doc~4}}}\hfill{\tiny\texttt{score:~10.996}}\ {\tiny\textcolor{gray}{\{'relevant': \{'53573702-4140-4812-8658-ae10a54edfc8': \{'value': 'yes', 'timestamp': 1723040808\}, '9dcd142c-415e-4161-a82f-c8b14bfb36ac': \{'value': 'no', 'timestamp': 1726570250\}\}\}}}\\[3pt]
{\small\textit{Discovering Alien Oceans: Magnetism – NASA's Europa Clipper}}\\[2pt]
{\tiny\url{https://europa.nasa.gov/resources/446/discovering-alien-oceans-magnetism/}}\\[4pt]
{\footnotesize Discovering Alien Oceans: Magnetism – NASA's Europa Clipper
An identical copy of Europa Clipper's propulsion module and electronics vault is lowered by crane in a static test tower.

Europa Clipper's Developmental Test Model

Europa Clipper’s 28-foot (8.5-meter) magnetometer boom is unfurled at NASA’s Jet Propulsion Laboratory.

Europa Clipper Team Deploys Magnetometer Boom

This video shows a test deployment of Europa Clipper’s very high frequency (VHF) radar antenna.

Europa Clipper REASON VHF Radar Antenna Deployment

This color composite view combines violet, green, and infrared images of Jupiter's intriguing moon, Europa, for a view of the moon in natural color (left) and in enhanced color designed to bring ou...

Europa Global Views in Natural and Enhanced Colors

Workers at NASA’s Jet Propulsion Laboratory in Southern California deploy Europa Clipper’s magnetometer boom.

Workers Deploy the Europa Clipper Magnetometer Boom

This view of Jupiter's icy moon Europa shows a region shaped like a mitten that has a texture similar to the matrix of chaotic terrain, which is seen in medium and high resolution images of numerou...

 Mitten Shaped Region of Chaotic Terrain on Europa 

This complex area on the side of Europa which faces away from Jupiter shows several types of features which are formed by disruptions of Europa's icy crust.

Scrambled Ice

Europa Clipper's high-gain antenna is installed on Aug. 14, 2023.

Powerful Antenna Installed On Europa Clipper

This image of Europa, an icy satellite of Jupiter about the size of the Earth's Moon, was obtained from a range of 7,415 miles (11,933 kilometers) by the Galileo spacecraft during its fourth orbit ...

Macula on Europa 

This enhanced color image from NASa's Galileo mission shows the region surrounding the young impact crater Pwyll on Jupiter's moon Europa.

Pwyll Crater

This artist's rendering shows NASA's Europa Clipper spacecraft, which is aiming for launch readiness by 2024.

2021 Europa Clipper Spacecraft (Artist's Concept)

This animation segment shows NASA’s Europa Clipper flying by Mars during its Mars gravity assist. 

Mission Animation: Mars Gravity Assist}
\end{docenv}
\begin{docenv}
{\footnotesize\textcolor{doccolor}{\textbf{Doc~5}}}\hfill{\tiny\texttt{score:~14.094}}\ {\tiny\textcolor{gray}{\{'relevant': \{'9dcd142c-415e-4161-a82f-c8b14bfb36ac': \{'value': 'yes', 'timestamp': 1726570311\}\}\}}}\\[3pt]
{\small\textit{Meet Europa Clipper | Spacecraft – NASA's Europa Clipper}}\\[2pt]
{\tiny\url{https://europa.nasa.gov/spacecraft}}\\[4pt]
{\footnotesize Meet Europa Clipper | Spacecraft – NASA's Europa Clipper
Europa Clipper’s propulsion module is an aluminum cylinder 10 feet (3 meters) long and 5 feet (1.5 meters) wide. It holds the spacecraft’s 24 engines and nearly 6,000 pounds (2,750 kilograms) of fuel in tanks, as well as the spacecraft’s helium pressurant tanks. 
Vault

Europa Clipper's electronics are enclosed in a vault with walls made of 1/3-inch thick (9.2-mm) sheets of aluminum-zinc alloy to protect the electronics from Jupiter’s intense radiation. These electronics include computers (or the “brains” of the spacecraft), flight software, and more. The vault is bolted to the spacecraft’s propulsion module and has wires to provide communications and control across the entire spacecraft.
Power

Europa Clipper’s “wings” - its solar arrays - span more than 100 feet (30.5 meters). Each wing is approximately 46.5 feet (14.2 meters) long and about 13.5 feet (4.1 meters) high. The huge arrays are needed to collect sunlight to power the spacecraft while it operates at Jupiter, which is more than five times as far from the Sun as Earth is.
Communications 

Europa Clipper’s 10-foot-diameter (3-meter-diameter) high-gain antenna (HGA) and other smaller radio antennas will be used to receive commands from Earth through NASA’s Deep Space Network. The antennas also will send data back to engineers and scientists. Europa Clipper’s antennas also will be used for gravity and radio science investigations, allowing scientists to study Europa’s gravity as the moon orbits Jupiter.
Guidance, Navigation, \& Control

Europa Clipper will use sensors, cameras, and special hardware to determine and control its position in space. Europa Clipper’s two stellar reference units, informally called “star trackers,” are small cameras with a computer that carries a catalog of stars with known positions. 
As the cameras observes the starry sky, the computer matches the pattern of stars in the images to its star catalog and determines the spacecraft's orientation in space. Europa Clipper will also use four reaction wheels to stay pointed in the right direction.}
\end{docenv}
\begin{refanswerenv}
{\footnotesize\textit{\textcolor{gray}{\textbf{Reference answer:}\enspace Yes, the Europa Clipper spacecraft has an high-gain antenna. It is used for communication with Earth and sending data back to NASA. Workers were seen installing it on the spacecraft on August 14, 2023.}}}
\end{refanswerenv}
\vspace{0.5em}
\noindent{\small\textcolor{usercolor}{\textbf{$\triangleright$\enspace User --- Turn~7}}}\hfill{\tiny\textcolor{gray}{Factoid $\cdot$ Follow-up $\cdot$ ANSWERABLE}}
\begin{userbubble}
What is the speed at which Clipper flies?
\end{userbubble}
\noindent\hfill{\small\textcolor{agentcolor}{\textbf{Agent\enspace$\triangleleft$}}}
\begin{agentbubble}
The propulsion subsystem of the Europa Clipper will carry fuel and oxidizer to change its velocity by about 3,600 miles per hour (5,800 kilometers per hour). While this is a significant and crucial capability, it represents only about five percent of the velocity change needed for the mission. The majority of the velocity change required to adjust trajectories once in Jupiter's orbit will come from planned gravitational flybys of Jovian moons. The propulsion subsystem can be thought of as the "gas tank" of the Europa Clipper.
\end{agentbubble}
\noindent\hfill{\footnotesize\textcolor{doccolor}{\textbf{$\hookrightarrow$\enspace Retrieved documents~(5)}}}
\begin{docenv}
{\footnotesize\textcolor{doccolor}{\textbf{Doc~1}}}\hfill{\tiny\texttt{score:~47.200}}\ {\tiny\textcolor{gray}{\{'relevant': \{'53573702-4140-4812-8658-ae10a54edfc8': \{'value': 'no', 'timestamp': 1723041676\}\}\}}}\\[3pt]
{\small\textit{Science | Mission – NASA's Europa Clipper}}\\[2pt]
{\tiny\url{https://europa.nasa.gov/mission/science/}}\\[4pt]
{\footnotesize Science | Mission – NASA's Europa Clipper
Science | Mission – NASA's Europa Clipper

NASA

Europa Clipper

Skip Navigation

menu

close modal

Europa Clipper Mission

Science

Mission

MENU

Overview

Science

Timeline

The Team

Overview

History

FAQ

Key Fact

What is Europa Clipper's main science goal? 

Europa Clipper’s main science goal is to determine whether there are places below the surface of Jupiter’s icy moon, Europa, that could support life.

OverviewOverview

Deep beneath the crust of Jupiter’s frozen moon Europa lies a massive liquid water ocean. Exploring this ocean world with our Europa Clipper spacecraft could provide new clues in our search for life beyond Earth. Scheduled to launch in October 2024, Europa Clipper’s main science goal is to determine whether there are places below the surface of Jupiter’s icy moon, Europa, that could support life. Credit: NASA

Jupiter’s moon Europa shows strong evidence for an ocean of liquid water beneath its icy crust. Beyond Earth, Europa is considered one of the most promising currently habitable environments in our solar system. 
Below Europa’s icy surface, evidence suggests there is a global ocean with more water than all of Earth’s oceans combined. Europa could have all the “ingredients” needed for life as we know it: 

Water: Twice as much as Earth’s oceans 

Organics: Essential chemical building blocks from a variety of sources

Energy: Chemical energy sources from the surface and the sea floor

Stability: Conditions remaining similar for 4 billion years

NASA's Europa Clipper spacecraft will perform approximately 50 close flybys of the moon, gathering detailed measurements to investigate whether the moon could have conditions suitable for life. Europa Clipper is not a life detection mission – its main science goal is to determine whether there are places below Europa’s surface that could support life.

Europa Up Close

Take an in-depth look at Jupiter's icy moon and Europa Clipper's plans for investigation.

Learn more ›

Science ObjectivesScience Objectives
Europa Clipper has three main science objectives: 

Determine the Thickness of Europa’s Icy Shell \& How the Ocean Interacts With the Surface}
\end{docenv}
\begin{docenv}
{\footnotesize\textcolor{doccolor}{\textbf{Doc~2}}}\hfill{\tiny\texttt{score:~45.433}}\ {\tiny\textcolor{gray}{\{'relevant': \{'53573702-4140-4812-8658-ae10a54edfc8': \{'value': 'no', 'timestamp': 1723041700\}\}\}}}\\[3pt]
{\small\textit{Europa Clipper - NASA Science}}\\[2pt]
{\tiny\url{https://science.nasa.gov/mission/europa-clipper/}}\\[4pt]
{\footnotesize Europa Clipper - NASA Science
Scientists predict Europa has a salty ocean beneath its icy crust that could hold the building blocks necessary to sustain life.TypeOrbiterLaunchOct. 10, 2024TargetEuropaObjectiveDetermine if Europa has conditions suitable to support lifeMission OverviewSpacecraftWhy Europa?What will Europa Clipper do?Europa Clipper’s main science goal is to determine whether there are places below the surface of Jupiter’s icy moon, Europa, that could support life.The mission’s three main science objectives are to understand the nature of the ice shell and the ocean beneath it, along with the moon’s composition and geology. The mission’s detailed exploration of Europa will help scientists better understand the astrobiological potential for habitable worlds beyond our planet.Read MoreAn artist's concept of NASA’s Europa Clipper spacecraft.NASA/JPL-CaltechWhat is NASA’s Europa Clipper spacecraft?Europa Clipper is a robotic solar-powered spacecraft built to conduct the first detailed investigations of Jupiter's icy moon Europa.With its solar arrays deployed, Europa Clipper spans more than 100 feet (about 30 meters) – about the length of a basketball court. The main body of the spacecraft consists of its avionics vault, radiofrequency module, and propulsion module.Read MoreWorkers use a crane to lift Europa Clipper's high-gain antenna as they prepare to install it on the spacecraft on Aug. 14, 2023. NASA/JPL-CaltechWhy Europa?The search for life beyond Earth is one of NASA’s primary objectives. If humans are to truly understand our place in the universe, we must learn whether our planet is the only place where life exists. So the search is on! There is strong evidence Jupiter's moon Europa has a saltwater ocean that may be one of the best places to look for environments where life could exist beyond Earth.Read MoreThis view of Jupiter’s icy moon Europa was captured by the JunoCam imager aboard NASA’s Juno spacecraft during the mission’s close flyby on Sept. 29, 2022.}
\end{docenv}
\begin{docenv}
{\footnotesize\textcolor{doccolor}{\textbf{Doc~3}}}\hfill{\tiny\texttt{score:~44.823}}\ {\tiny\textcolor{gray}{\{'relevant': \{'53573702-4140-4812-8658-ae10a54edfc8': \{'value': 'no', 'timestamp': 1723041725\}\}\}}}\\[3pt]
{\small\textit{Europa Clipper High Gain Antenna Undergoes Testing at Langley – NASA's Europa Clipper}}\\[2pt]
{\tiny\url{https://europa.nasa.gov/news/23/europa-clipper-high-gain-antenna-undergoes-testing-at-langley/}}\\[4pt]
{\footnotesize Europa Clipper High Gain Antenna Undergoes Testing at Langley – NASA's Europa Clipper
Jupiter’s moon Europa and its global ocean may currently have conditions suitable for life. Scientists are studying processes on the icy surface as they prepare to explore.

Surface of Jupiter's Moon Europa Churned by Small Impacts

Jupiter’s icy moon Europa may be the most promising place in the solar system to find present-day environments suitable for life beyond Earth. 

Europa: A World of Ice, With Potential for Life

Jupiter’s moon Europa has an icy crust covering a vast, global ocean. The rocky layer underneath may be hot enough to melt, leading to undersea volcanoes.

Europa's Interior May Be Hot Enough to Fuel Seafloor Volcanoes

Jupiter’s moon Europa may have the potential to harbor life. The spacecraft will use multiple flybys of the moon to investigate the habitability of this ocean world.

NASA's Europa Clipper Builds Hardware, Moves Toward Assembly

In an unprecedented year, the Europa Clipper project continues on the path toward unprecedented science. 

On the Path Towards Unprecedented Science

Scientists have theorized on the origin of the water plumes possibly erupting from Jupiter's moon Europa. Recent research adds a new potential source to the mix.

Potential Plumes on Europa Could Come From Water in the Crust

New lab experiments re-create the environment of Europa and find that the icy moon shines, even on its nightside. The effect is more than just a cool visual.

Europa Glows: Radiation Does a Bright Number on Jupiter's Moon

The surface of Jupiter's moon Europa features a widely varied landscape, including ridges, bands, small rounded domes and disrupted spaces that geologists call "chaos terrain." Three newly reprocessed images, taken by NASA's Galileo spacecraft in the late 1990s, reveal details in diverse surface features on Europa.

Newly Reprocessed Images of Europa Show 'Chaos Terrain' in Crisp Detail

An icy ocean world in our solar system that could tell us more about the potential for life on other worlds is coming into focus with confirmation of the Europa Clipper mission's next phase. 

Mission to Jupiter's Icy Moon Confirmed}
\end{docenv}
\begin{docenv}
{\footnotesize\textcolor{doccolor}{\textbf{Doc~4}}}\hfill{\tiny\texttt{score:~15.218}}\ {\tiny\textcolor{gray}{\{'relevant': \{'53573702-4140-4812-8658-ae10a54edfc8': \{'value': 'yes', 'timestamp': 1723047888\}\}\}}}\\[3pt]
{\small\textit{On the Path Towards Unprecedented Science – NASA's Europa Clipper}}\\[2pt]
{\tiny\url{https://europa.nasa.gov/news/28/on-the-path-towards-unprecedented-science/}}\\[4pt]
{\footnotesize On the Path Towards Unprecedented Science – NASA's Europa Clipper
It will carry fuel and oxidizer sufficient to change Europa Clipper’s velocity by about 3,600 miles per hour (5,800 kilometers per hour). While this is a large amount and a critically important capability, it represents about five percent of the velocity change needed to accomplish this mission. The bulk of the velocity change necessary to adjust trajectories once in Jupiter orbit will come from gravitational flybys of Jovian moons as planned. Given the overall size of the tanks, the propulsion subsystem is designed as part of a “propulsion module” that includes the solar arrays and radiation monitors.

Europa Clipper’s bowl-shaped high-gain antenna faces down in this image. This radio dish will allow ground controllers to send and receive commands and data between Earth and the spacecraft in Jupiter orbit — more than one million times farther from Earth than the International Space Station orbits. Credit: NASA

The telecom subsystem is the link between Europa Clipper and Earth. It provides command reception and science and engineering data downlink. It also provides interplanetary navigation data. Through the high gain antenna (HGA), Europa Clipper will downlink on the order of several gigabytes of science data per typical Earth antenna tracking session, enough data to fill a smartphone in about a week. The HGA also serves as a sunshade when Europa Clipper is still relatively close to the Sun, while en route to the Jupiter system. With respect to its gravity science role, scientists will use the radio signals exchanged between the telecom sub-system aboard the spacecraft and NASA's Deep Space Network to detect subtle variations in the gravitational attraction between Europa and Europa Clipper. The variations reflect details of Europa’s internal mass distribution, allowing us to learn about Europa’s interior structure.

The Europa Clipper mission intends to answer the big question: Is Europa habitable? This question encompasses hypotheses about the nature of water, chemistry, and energy exchange at Europa, which in turn leads to these additional major questions to be answered: How deep and salty is the ocean? How thick is the ice shell? How active is the ice shell? Are there plumes, and if so, what is in the plumes?}
\end{docenv}
\begin{docenv}
{\footnotesize\textcolor{doccolor}{\textbf{Doc~5}}}\hfill{\tiny\texttt{score:~14.187}}\ {\tiny\textcolor{gray}{\{'relevant': \{'53573702-4140-4812-8658-ae10a54edfc8': \{'value': 'yes', 'timestamp': 1723047967\}\}\}}}\\[3pt]
{\small\textit{On the Path Towards Unprecedented Science – NASA's Europa Clipper}}\\[2pt]
{\tiny\url{https://europa.nasa.gov/news/28/on-the-path-towards-unprecedented-science/}}\\[4pt]
{\footnotesize On the Path Towards Unprecedented Science – NASA's Europa Clipper
Shown here is an interactive rendering of the spacecraft. Fully deployed, the spacecraft will be more than 5 meters (15 feet) tall, 17 meters (50 feet) long, and will span about 30 meters (100 feet) in width — longer than a basketball court from tip to tip of its solar arrays. Credit: NASA/JPL-Caltech

Download Options

The solar panels and overall physical structure visible in the interactive model are aspects of spacecraft subsystems that comprise the flight system. The science instruments, attached directly to the spacecraft or through a boom, are referred to as the payload. Europa Clipper’s antennas fall into both of these categories; some antennas provide only a spacecraft (telecom) function, other antennas provide only a science function (for the radar), and yet others perform both telecom and science (gravity and radio science) roles.
The Europa Clipper flight system design incorporates lessons learned from prior missions; the scientific instruments are at least partly based on interplanetary mission predecessors. Informed by those experiences, the spacecraft and its components are designed to be appropriately robust to the Europa environment. Europa Clipper’s components also need to be mutually compatible with one another while being operated. Given that, the testing and performance verification is designed to reflect “flying conditions” as well as possible. Let’s look at two subsystems whose fabrication is underway.

NASA’s Goddard Space Flight Center (GSFC) engineers align the upper and lower cylinders of Europa Clipper’s propulsion module for a fit check. Credit: NASA/B. Lambert

The propulsion subsystem is essentially Europa Clipper’s “gas tank.” It will carry fuel and oxidizer sufficient to change Europa Clipper’s velocity by about 3,600 miles per hour (5,800 kilometers per hour). While this is a large amount and a critically important capability, it represents about five percent of the velocity change needed to accomplish this mission. The bulk of the velocity change necessary to adjust trajectories once in Jupiter orbit will come from gravitational flybys of Jovian moons as planned. Given the overall size of the tanks, the propulsion subsystem is designed as part of a “propulsion module” that includes the solar arrays and radiation monitors.}
\end{docenv}
\begin{refanswerenv}
{\footnotesize\textit{\textcolor{gray}{\textbf{Reference answer:}\enspace The speed at which the Europa Clipper spacecraft flies will vary depending on its trajectory and distance from Jupiter. However, during its closest flybys of Europa, it is expected to reach speeds of around 15,000 to 20,000 miles per hour (6.7 to 9 kilometers per second).}}}
\end{refanswerenv}
\vspace{0.5em}
\noindent{\small\textcolor{usercolor}{\textbf{$\triangleright$\enspace User --- Turn~8}}}\hfill{\tiny\textcolor{gray}{Factoid $\cdot$ Follow-up $\cdot$ ANSWERABLE}}
\begin{userbubble}
Does Clipper fly by Europe?
\end{userbubble}
\noindent\hfill{\small\textcolor{agentcolor}{\textbf{Agent\enspace$\triangleleft$}}}
\begin{agentbubble}
Yes, the main purpose of the Europa Clipper mission is to perform approximately 50 close flybys of Jupiter's moon Europa, gathering detailed measurements to investigate whether the moon could have conditions suitable for life.
\end{agentbubble}
\noindent\hfill{\footnotesize\textcolor{doccolor}{\textbf{$\hookrightarrow$\enspace Retrieved documents~(3)}}}
\begin{docenv}
{\footnotesize\textcolor{doccolor}{\textbf{Doc~1}}}\hfill{\tiny\texttt{score:~47.200}}\ {\tiny\textcolor{gray}{\{'relevant': \{'53573702-4140-4812-8658-ae10a54edfc8': \{'value': 'no', 'timestamp': 1723048315\}, '9dcd142c-415e-4161-a82f-c8b14bfb36ac': \{'value': 'yes', 'timestamp': 1726571105\}\}\}}}\\[3pt]
{\small\textit{Science | Mission – NASA's Europa Clipper}}\\[2pt]
{\tiny\url{https://europa.nasa.gov/mission/science/}}\\[4pt]
{\footnotesize Science | Mission – NASA's Europa Clipper
Science | Mission – NASA's Europa Clipper

NASA

Europa Clipper

Skip Navigation

menu

close modal

Europa Clipper Mission

Science

Mission

MENU

Overview

Science

Timeline

The Team

Overview

History

FAQ

Key Fact

What is Europa Clipper's main science goal? 

Europa Clipper’s main science goal is to determine whether there are places below the surface of Jupiter’s icy moon, Europa, that could support life.

OverviewOverview

Deep beneath the crust of Jupiter’s frozen moon Europa lies a massive liquid water ocean. Exploring this ocean world with our Europa Clipper spacecraft could provide new clues in our search for life beyond Earth. Scheduled to launch in October 2024, Europa Clipper’s main science goal is to determine whether there are places below the surface of Jupiter’s icy moon, Europa, that could support life. Credit: NASA

Jupiter’s moon Europa shows strong evidence for an ocean of liquid water beneath its icy crust. Beyond Earth, Europa is considered one of the most promising currently habitable environments in our solar system. 
Below Europa’s icy surface, evidence suggests there is a global ocean with more water than all of Earth’s oceans combined. Europa could have all the “ingredients” needed for life as we know it: 

Water: Twice as much as Earth’s oceans 

Organics: Essential chemical building blocks from a variety of sources

Energy: Chemical energy sources from the surface and the sea floor

Stability: Conditions remaining similar for 4 billion years

NASA's Europa Clipper spacecraft will perform approximately 50 close flybys of the moon, gathering detailed measurements to investigate whether the moon could have conditions suitable for life. Europa Clipper is not a life detection mission – its main science goal is to determine whether there are places below Europa’s surface that could support life.

Europa Up Close

Take an in-depth look at Jupiter's icy moon and Europa Clipper's plans for investigation.

Learn more ›

Science ObjectivesScience Objectives
Europa Clipper has three main science objectives: 

Determine the Thickness of Europa’s Icy Shell \& How the Ocean Interacts With the Surface}
\end{docenv}
\begin{docenv}
{\footnotesize\textcolor{doccolor}{\textbf{Doc~2}}}\hfill{\tiny\texttt{score:~45.433}}\ {\tiny\textcolor{gray}{\{'relevant': \{'53573702-4140-4812-8658-ae10a54edfc8': \{'value': 'no', 'timestamp': 1723048354\}\}\}}}\\[3pt]
{\small\textit{Europa Clipper - NASA Science}}\\[2pt]
{\tiny\url{https://science.nasa.gov/mission/europa-clipper/}}\\[4pt]
{\footnotesize Europa Clipper - NASA Science
Scientists predict Europa has a salty ocean beneath its icy crust that could hold the building blocks necessary to sustain life.TypeOrbiterLaunchOct. 10, 2024TargetEuropaObjectiveDetermine if Europa has conditions suitable to support lifeMission OverviewSpacecraftWhy Europa?What will Europa Clipper do?Europa Clipper’s main science goal is to determine whether there are places below the surface of Jupiter’s icy moon, Europa, that could support life.The mission’s three main science objectives are to understand the nature of the ice shell and the ocean beneath it, along with the moon’s composition and geology. The mission’s detailed exploration of Europa will help scientists better understand the astrobiological potential for habitable worlds beyond our planet.Read MoreAn artist's concept of NASA’s Europa Clipper spacecraft.NASA/JPL-CaltechWhat is NASA’s Europa Clipper spacecraft?Europa Clipper is a robotic solar-powered spacecraft built to conduct the first detailed investigations of Jupiter's icy moon Europa.With its solar arrays deployed, Europa Clipper spans more than 100 feet (about 30 meters) – about the length of a basketball court. The main body of the spacecraft consists of its avionics vault, radiofrequency module, and propulsion module.Read MoreWorkers use a crane to lift Europa Clipper's high-gain antenna as they prepare to install it on the spacecraft on Aug. 14, 2023. NASA/JPL-CaltechWhy Europa?The search for life beyond Earth is one of NASA’s primary objectives. If humans are to truly understand our place in the universe, we must learn whether our planet is the only place where life exists. So the search is on! There is strong evidence Jupiter's moon Europa has a saltwater ocean that may be one of the best places to look for environments where life could exist beyond Earth.Read MoreThis view of Jupiter’s icy moon Europa was captured by the JunoCam imager aboard NASA’s Juno spacecraft during the mission’s close flyby on Sept. 29, 2022.}
\end{docenv}
\begin{docenv}
{\footnotesize\textcolor{doccolor}{\textbf{Doc~3}}}\hfill{\tiny\texttt{score:~44.823}}\ {\tiny\textcolor{gray}{\{'relevant': \{'53573702-4140-4812-8658-ae10a54edfc8': \{'value': 'no', 'timestamp': 1723048435\}\}\}}}\\[3pt]
{\small\textit{Europa Clipper High Gain Antenna Undergoes Testing at Langley – NASA's Europa Clipper}}\\[2pt]
{\tiny\url{https://europa.nasa.gov/news/23/europa-clipper-high-gain-antenna-undergoes-testing-at-langley/}}\\[4pt]
{\footnotesize Europa Clipper High Gain Antenna Undergoes Testing at Langley – NASA's Europa Clipper
Jupiter’s moon Europa and its global ocean may currently have conditions suitable for life. Scientists are studying processes on the icy surface as they prepare to explore.

Surface of Jupiter's Moon Europa Churned by Small Impacts

Jupiter’s icy moon Europa may be the most promising place in the solar system to find present-day environments suitable for life beyond Earth. 

Europa: A World of Ice, With Potential for Life

Jupiter’s moon Europa has an icy crust covering a vast, global ocean. The rocky layer underneath may be hot enough to melt, leading to undersea volcanoes.

Europa's Interior May Be Hot Enough to Fuel Seafloor Volcanoes

Jupiter’s moon Europa may have the potential to harbor life. The spacecraft will use multiple flybys of the moon to investigate the habitability of this ocean world.

NASA's Europa Clipper Builds Hardware, Moves Toward Assembly

In an unprecedented year, the Europa Clipper project continues on the path toward unprecedented science. 

On the Path Towards Unprecedented Science

Scientists have theorized on the origin of the water plumes possibly erupting from Jupiter's moon Europa. Recent research adds a new potential source to the mix.

Potential Plumes on Europa Could Come From Water in the Crust

New lab experiments re-create the environment of Europa and find that the icy moon shines, even on its nightside. The effect is more than just a cool visual.

Europa Glows: Radiation Does a Bright Number on Jupiter's Moon

The surface of Jupiter's moon Europa features a widely varied landscape, including ridges, bands, small rounded domes and disrupted spaces that geologists call "chaos terrain." Three newly reprocessed images, taken by NASA's Galileo spacecraft in the late 1990s, reveal details in diverse surface features on Europa.

Newly Reprocessed Images of Europa Show 'Chaos Terrain' in Crisp Detail

An icy ocean world in our solar system that could tell us more about the potential for life on other worlds is coming into focus with confirmation of the Europa Clipper mission's next phase. 

Mission to Jupiter's Icy Moon Confirmed}
\end{docenv}
\end{document}
